%==========================================================================
% RESUMO (Português)
%==========================================================================
\chapter*{Resumo}
\addcontentsline{toc}{chapter}{Resumo}

A viticultura é uma das atividades agrícolas com maior relevância económica em
Portugal e na Europa. A identificação precisa dos gomos nas videiras é uma
etapa fundamental para a estimativa da produção e para a realização da poda
mecanizada — tarefa que, se automatizada, pode reduzir significativamente os
custos operacionais e a dependência de mão de obra especializada.

O presente relatório descreve o trabalho desenvolvido durante o estágio
curricular realizado no \textbf{Instituto Politécnico de Bragança (IPB)}, no
âmbito do \textbf{Centro de Investigação em Digitalização e Robótica
Inteligente (CEDRI)}, entre março e julho de 2026. O tema central é a
\textbf{deteção e classificação automática de gomos em videiras} utilizando
técnicas de visão computacional e aprendizagem profunda (\textit{deep
learning}).

Durante o estágio foram desenvolvidos três modelos progressivos. O primeiro
baseou-se no dataset público \textbf{Segmentation-Vineyard-2}~\cite{vinya2024segvineyard}
(origem espanhola), que continha anotações de \textit{sarmento} e
\textit{tronco} — tendo sido adicionadas manualmente as anotações de
\textit{gomo} pelo estagiário. Os modelos seguintes incorporaram imagens
retiradas de vídeos captados nas vinhas de Palmela, em condições próximas
de um sistema robótico real de corte.
Estas imagens apresentam elevada complexidade visual, com oclusões, planos
de fundo densos (outras videiras, árvores, postes e arames), múltiplos
elementos em cena e variações de iluminação. Para lidar com essa complexidade,
foram anotadas quatro classes de objetos: \textit{gomos}, \textit{sarmentos},
\textit{troncos} e \textit{objetos} (arames, postes e outros elementos).

O modelo de deteção selecionado foi o \textbf{YOLO26 Nano}~\cite{hidayatullah2026yolo26},
pela sua eficiência computacional, inferência \textit{NMS-free} e adequação a sistemas
embebidos. O processo de anotação e treino foi realizado na plataforma
\textbf{Roboflow}. O Modelo~2 (400 imagens base, 772 após \textit{data
augmentation}) obteve mAP@50 de 48,8\%. O Modelo~3 (938 imagens após
augmentation) obteve mAP@50 de 43,7\% e Precisão de 57,3\%, com os gomos
a surgir nas inferências a limiares de confiança de 16--26\%.

Os resultados obtidos evidenciam que a deteção de gomos em cenas reais
é um problema exigente, com a complexidade visual a requerer volumes de
dados superiores ao que pode ser alcançado num único semestre. As lições
aprendidas e os resultados dos três modelos constituem uma base sólida
para trabalhos futuros de maior escala.

\vspace{1em}
\noindent\textbf{Palavras-chave:} visão computacional, \textit{deep learning},
YOLO26, deteção de objetos, viticultura, gomos, poda mecanizada, Roboflow.

\clearpage

%==========================================================================
% ABSTRACT (Inglês)
%==========================================================================
\chapter*{Abstract}
\addcontentsline{toc}{chapter}{Abstract}

Viticulture is one of the most economically significant agricultural activities
in Portugal and across Europe. The accurate identification of buds on
grapevines is a critical step for yield estimation and mechanised pruning —
a task that, if automated, can substantially reduce operational costs and
reliance on skilled labour.

This report describes the work carried out during a curricular internship at
the \textbf{Polytechnic Institute of Bragança (IPB)}, within the \textbf{Research
Centre in Digitalization and Intelligent Robotics (CEDRI)}, from March to
July 2026. The central theme is the \textbf{automatic detection and
classification of buds on grapevines} using computer vision and deep learning
techniques.

Three progressive models were developed during the internship. The first was
based on the publicly available \textit{Segmentation-Vineyard-2} Spanish
dataset, to which bud annotations were manually added. Subsequent models
incorporated frames extracted from videos captured in the Palmela vineyards,
under conditions close to those of a real robotic pruning system.
These images exhibit high visual complexity, including occlusions, dense
backgrounds (other vines, trees, poles and wires), multiple elements in the
scene and lighting variations. Four object classes were annotated:
\textit{buds}, \textit{canes (sarmentos)}, \textit{trunks}
and \textit{objects} (wires, poles and other structural elements).

The selected detection model was \textbf{YOLO26 Nano}~\cite{hidayatullah2026yolo26},
chosen for its computational efficiency, NMS-free inference and suitability
for embedded systems. Annotation and training were performed on the
\textbf{Roboflow} platform. Model~2 (400 base images, expanded to 772 by
data augmentation) achieved a mAP@50 of 48.8\%. Model~3 (938 images after
augmentation) achieved mAP@50 of 43.7\% and Precision of 57.3\%, with buds
appearing in inferences at confidence thresholds of 16--26\%.

Results demonstrate that bud detection in real vineyard scenes is a
challenging problem. The three models developed establish a solid methodological
foundation and provide a realistic assessment of the data volumes required
for a production-grade system.

\vspace{1em}
\noindent\textbf{Keywords:} computer vision, deep learning, YOLO26, object
detection, viticulture, buds, mechanised pruning, Roboflow.

\clearpage
